Un accident de la route peut avoir des conséquences très différentes pour les personnes
impliquées. Certains usagers sont indemnes, tandis que
d'autres sont blessés, hospitalisés ou décèdent. Comprendre ces différences conduit
à examiner les caractéristiques des personnes, les véhicules et les circonstances
de l'accident.

Ce projet s'appuie sur les données du Bulletin d'analyse des accidents corporels de la
circulation, ou BAAC, publiées sur
data.gouv.fr\footnote{\href{https://www.data.gouv.fr/datasets/bases-de-donnees-annuelles-des-accidents-corporels-de-la-circulation-routiere-annees-de-2005-a-2024}{ONISR, bases de données annuelles des accidents corporels de la circulation, data.gouv.fr.}}.
L'étude couvre la France métropolitaine, Corse comprise, de 2019 à 2024. Les fichiers
décrivent les accidents corporels enregistrés par les forces de l'ordre, leurs lieux,
les véhicules et les usagers concernés. Notre unité d'analyse est la participation d'un
usager à un accident. Les probabilités étudiées portent donc sur la gravité parmi les
personnes présentes dans ces enregistrements.

La problématique retenue est la suivante~:
\begin{quote}
\itshape
Quelles caractéristiques sont associées au décès ou à l'hospitalisation des usagers
impliqués dans les accidents corporels enregistrés en métropole, et dans quelle mesure
peut-on estimer la probabilité de cette issue sur une année ultérieure~?
\end{quote}

La première difficulté consiste à construire une base cohérente à partir des quatre
fichiers de chaque année. Les informations décrivent des niveaux différents, les
identifiants nécessitent une harmonisation et plusieurs lieux peuvent être associés
au même accident. Une jointure mal définie pourrait ainsi compter plusieurs fois les
mêmes usagers. Les codes inconnus, les valeurs atypiques et les informations manquantes
demandent également des règles explicites. La préparation occupe donc une place centrale~:
ses décisions doivent être justifiées et leurs effets vérifiés.

L'analyse poursuit ensuite deux objectifs complémentaires. Le premier est descriptif
et statistique~: caractériser la population, comparer les proportions de décès ou
d'hospitalisations selon les groupes et étudier l'incertitude de ces écarts, en conservant
le regroupement des usagers par accident. Le second est prédictif~: évaluer si plusieurs
caractéristiques considérées ensemble permettent d'estimer ces probabilités mieux
qu'une référence simple. Ce volet repose sur un apprentissage
en 2019-2022, une validation en 2023 et une évaluation finale en 2024. Il porte sur
des accidents déjà enregistrés~; la portée des prédictions dépend des informations
retenues pour l'estimation et de leur instant de disponibilité.

Les résultats sont interprétés au regard des limites de collecte, des informations
inconnues et de la distinction entre association statistique et causalité. Les limites
du regroupement «~décès ou hospitalisation~» sont également discutées.

La première partie du rapport présente le cadrage et le cahier des charges, les sources
BAAC, leur audit, la préparation, puis les analyses exploratoires et statistiques.
Elle explique les difficultés rencontrées, les solutions retenues et les contrôles
réalisés. La seconde présente le protocole prédictif, la référence constante, la régression
logistique et leur évaluation. Une conclusion générale et un glossaire complètent
le rapport.
